\section{A common purification of effects}
\label{sec:operators}

Conditioning on answers replaces the shared state by local positive
operators applied to that state. We need purifications of these operators
that are close whenever an operator entropy gap is small. For positive
definite operators, this gap is an average of Petz's operator Bregman
divergence~\cite{petz2007}; related resolvent estimates appear in
Kim~\cite{kim2014}. The following construction provides all purifications
in a common space and will let us retain the original answer measurements.
The purification formulation and its Jensen estimate are given in
\cite[Chapter~6, Definition~4.2, Lemma~4.3, and Appendix~A.4]{tenadvances2026};
we include the proof, including its finite-dimensional reduction.

For a finite-dimensional complex Hilbert space $\mathcal H$, an
\emph{effect} is a self-adjoint operator $T$ with $0\preceq T\preceq I$.
Put $f(t)=t\log t$ for $t>0$ and $f(0)=0$, and define
\[
 \mathcal F_T:\mathcal H\longrightarrow
 L^2((0,\infty),ds;\mathcal H),\qquad
 (\mathcal F_T h)(s)=T(T+sI)^{-1}h.
\]
The inverse in this expression exists for every $s>0$, including when
$T$ has a kernel.

\begin{lemma}[Resolvent purification]
\label{lem:resolvent}
For every effect $T$ on $\mathcal H$,
$\mathcal F_T^*\mathcal F_T=T$.  If $A_1,\ldots,A_N$ are effects,
$q_1,\ldots,q_N>0$ sum to $1$, and $B=\sum_jq_jA_j$, then
\begin{equation}
 \sum_j q_j(\mathcal F_{A_j}-\mathcal F_B)^*
                 (\mathcal F_{A_j}-\mathcal F_B)
 \preceq \sum_jq_j f(A_j)-f(B).
 \label{eq:resolvent-jensen}
\end{equation}
For any finite collection $\mathscr T$ of effects, all maps
$\mathcal F_T$, $T\in\mathscr T$, can instead take values in one
finite-dimensional space, preserving every product
$\mathcal F_S^*\mathcal F_T$, $S,T\in\mathscr T$.
\end{lemma}

\begin{proof}
The scalar identity
\[
 \int_0^\infty\frac{t^2}{(t+s)^2}\,ds=t\qquad(t\ge0)
\]
and the spectral theorem give the first assertion.

To prove \eqref{eq:resolvent-jensen}, write $D_j=A_j-B$.
If $h\in\ker B$, positivity and $q_j>0$ imply $A_jh=0$ for
all $j$.  This common kernel reduces every operator in question,
and both sides of \eqref{eq:resolvent-jensen} vanish there.
We may therefore restrict to $\operatorname{supp}B$, where $B>0$.
The scalar integral
\[
 f(t)=\int_0^\infty
       \left(\frac{t}{1+s}-\frac{t}{t+s}\right)ds
 \qquad(t\ge0)
\]
is convergent also at $t=0$, with the integrand then zero.  Functional
calculus and $\sum_jq_jD_j=0$ give
\begin{align}
 \sum_jq_j f(A_j)-f(B)
 &=\int_0^\infty s\sum_jq_j
       \big((A_j+sI)^{-1}-(B+sI)^{-1}\big)\,ds \notag\\
 &=\int_0^\infty s\sum_jq_j(B+sI)^{-1}D_j
       (A_j+sI)^{-1}D_j(B+sI)^{-1}\,ds.
 \label{eq:resolvent-gap}
\end{align}
For the second equality use the exact identity
\begin{align*}
 (A_j+sI)^{-1}
 ={}&(B+sI)^{-1}-(B+sI)^{-1}D_j(B+sI)^{-1}\\
 &+(B+sI)^{-1}D_j(A_j+sI)^{-1}D_j(B+sI)^{-1},
\end{align*}
then average over $j$.  The final integrand in
\eqref{eq:resolvent-gap} is bounded near $s=0$, since $B>0$ and
$s(A_j+sI)^{-1}\preceq I$, and is $O(s^{-2})$ as
$s\to\infty$.  Thus these are norm-convergent operator integrals.

The resolvent identity also gives
\[
 A_j(A_j+sI)^{-1}-B(B+sI)^{-1}
   =s(A_j+sI)^{-1}D_j(B+sI)^{-1}.
\]
Consequently the left side of \eqref{eq:resolvent-jensen} has integrand
\[
 s^2\sum_jq_j(B+sI)^{-1}D_j
       (A_j+sI)^{-2}D_j(B+sI)^{-1}.
\]
Since
$s(A_j+sI)^{-2}\preceq(A_j+sI)^{-1}$, congruence by
$D_j(B+sI)^{-1}$ bounds this integrand by the final integrand in
\eqref{eq:resolvent-gap}.  Integration proves the inequality.

Finally, the span
\[
 \mathcal K=\operatorname{span}\{\mathcal F_T h:
              T\in\mathscr T,\ h\in\mathcal H\}
\]
has dimension at most $|\mathscr T|\dim\mathcal H$.
Viewing every map as taking values in this same subspace preserves all
inner products of their images, hence all the stated products.
An unused one-dimensional summand may be added if a nonzero codomain
is needed.
\end{proof}

The common codomain is essential: compressing different operators
separately would not in general preserve their distances.  In our
application the number of effects is finite, although their number and
the resulting dimension need not have a uniform bound.

The transport of answer effects below follows the ideal-state extraction
in \cite[Chapter~6, Lemma~3.2 and Appendix~A.1]{tenadvances2026}.

\begin{lemma}[Transporting answer measurements]
\label{lem:measurement-transport}
Let $\mathcal H,\mathcal K$ be finite-dimensional complex Hilbert spaces,
let $T$ be an effect on $\mathcal H$, and let
$F:\mathcal H\to\mathcal K$ satisfy $F^*F=T$.
Suppose that $\mathcal A$ is a nonempty finite set and that positive
operators $(G_a)_{a\in\mathcal A}$ satisfy $\sum_aG_a=T$.
There is a POVM $(\widetilde G_a)_{a\in\mathcal A}$ on $\mathcal K$
with
\begin{equation}
 F^*\widetilde G_aF=G_a\qquad(a\in\mathcal A).
 \label{eq:measurement-pullback}
\end{equation}
The POVM depends only on $F$ and the operators $G_a$.
\end{lemma}

\begin{proof}
All inverses below act on the support of $T$ and are zero on its
kernel.  Define $V_T=FT^{-1/2}$.  Then
$V_T^*V_T=\Pi_T$, the orthogonal projection onto
$\operatorname{supp}T$, so $V_TV_T^*$ is an orthogonal projection
on $\mathcal K$.  Each $G_a$ is supported on
$\operatorname{supp}T$, because $0\preceq G_a\preceq T$.
Fix $a_*\in\mathcal A$ and put
\[
 \widetilde G_a
 =V_TT^{-1/2}G_aT^{-1/2}V_T^*
   +\mathbf1_{\{a=a_*\}}(I_{\mathcal K}-V_TV_T^*).
\]
These operators are positive and sum to $I_{\mathcal K}$.
Since $F=V_TT^{1/2}$, their pullbacks are $G_a$, as claimed.
The same formula applies when $T=0$: all $G_a$ and $V_T$ then vanish,
and the POVM always outputs $a_*$.
\end{proof}

In particular, if $\Psi=(\mathcal F_T\otimes\mathcal F_S)\psi$
and $\sum_aG_a=T$, $\sum_bH_b=S$, the transported measurements satisfy
\begin{equation}
 \langle\Psi|\widetilde G_a\otimes\widetilde H_b|\Psi\rangle
   =\langle\psi|G_a\otimes H_b|\psi\rangle.
 \label{eq:joint-measurement-pullback}
\end{equation}
Thus the common purification preserves the exact answer probabilities,
not just the norms of the conditioned vectors.
